\section{Introduction}
\label{sec:intro}

Industrial fault-diagnosis research increasingly claims generalization:
models are praised for surviving \emph{operating-condition shift}---a bearing
trained at one load or on one test rig and evaluated on another. The claim is
only as good as the experiment that produces it, and those experiments are
fragile in specific, recurring ways. Vibration records are typically cut into
overlapping windows; unless windows of one recording are kept in a single
partition, near-duplicate samples straddle the train/test boundary and inflate
scores. Normalization statistics fitted over the whole dataset quietly fold
held-out conditions into training. Model selection peeks at the test
condition. A seed that crashed is dropped, leaving three reported runs of an
experiment registered with four. And when a result is challenged, the chain
from numbers back to data files, code, and configuration no longer exists.
Recent standardization work in the domain makes exactly this complaint: ad hoc
partitions and evaluation setups make cross-condition results incomparable
across papers~\cite{zhao2024dgbenchmark,ieeeaccess2025standard}.

\begin{figure}[t]
\centering
% Replace with your artwork: \includegraphics[width=\linewidth]{figures/figure1.pdf}
\fbox{\parbox[c][5.2cm][c]{0.96\linewidth}{\centering artwork placeholder}}
\caption{The same experiment under two regimes, measured on the real bearing
archive with identical model, policy, and seeds (audit policy; artefact:
\texttt{results/verification/leakage-motivated-example-2026-11-07.json}).
Left: a conventional window-level random split lets windows of one recording
straddle the train/test boundary (all 36 recordings shared across partitions
in every seed) and inflates mean test Macro-F1 to $0.449$ against $0.293$
under a recording-grouped split with zero shared recordings. Right: rather
than executing such an experiment, EdgeFault-Bench rejects the identical
artefact at the \textsc{TaskAudited} transition with a machine-readable
invariant violation (\emph{recording R spans partitions train/test}, INV-2);
nothing partial survives.}
\label{fig:motivated}
\end{figure}

A natural first reaction is to standardize the benchmark: fix the datasets,
the splits, and the metrics, as domain-generalization research did for vision
(DomainBed~\cite{gulrajani2021domainbed}) and as fault-diagnosis benchmark
studies have begun to do~\cite{zhao2024dgbenchmark,dgbenchmarkrepo}. This
standardizes \emph{what} is evaluated. It does not, however, guarantee
\emph{that the evaluation is valid}: leakage controls, test-set isolation,
seed completeness, and provenance remain conventions living in scripts, to be
re-implemented---and re-broken---by every extension. The complementary
reaction is infrastructure: experiment trackers such as MLflow~\cite{mlflow}
record parameters, code versions, and artefacts. But trackers are
record-keepers; they enforce nothing about partition exclusivity or seed
completeness, and integrity remains voluntary and after the fact.

This paper takes a third position: \emph{the benchmark itself should be a
software system under formal specification}, in which experiment integrity is
not a convention but a property enforced by architecture and checked by
tools. Concretely, we ask whether the validity of an industrial ML experiment
can be guaranteed by a small set of kernel-enforced invariants, whether
third-party datasets and models can be integrated without any ability to
violate those invariants, and whether the guarantee itself can be verified
rather than asserted.

Making that position real raises three challenges. First,
\emph{heterogeneity against integrity}: real datasets span file formats,
sampling rates (sometimes unstated by the source), sensor modalities, and
condition taxonomies, while model code spans classical estimators and deep
frameworks; every extension point is an opportunity to break an invariant.
Second, \emph{enforcement must be verifiable, not declared}: invariants that
live only in documentation are indistinguishable from absent ones; the guards
must be explicit, machine-checkable, and demonstrably effective against
injected violations. Third, \emph{validity must survive scale and
distribution}: hashing, streaming, and result bundles must work on
multi-gigabyte archives, and the framework's own cost must be measured and
acceptable.

EdgeFault-Bench answers these with a four-module design. A
\emph{data-and-protocol} module fixes canonical recording and window identity
and expresses tasks---including combined-condition and
leave-one-device-out designs---as versioned manifests rather than code. An
\emph{extensible execution} module streams signals from checksum-verified
sources through dataset, preprocessing, and model plugins; training plugins
receive partition-restricted views and are architecturally unable to read the
test partition, which is materialized only inside the evaluation step. A
\emph{formal integrity} module specifies the experiment lifecycle as a guarded
state machine over six invariants, checked exhaustively with
TLA\textsuperscript{+}/TLC~\cite{lamport2002specifying} at the design level
and continuously by property-based tests~\cite{hypothesis} and a
fault-injection catalogue at the implementation level. A
\emph{reproducibility} module emits result bundles whose inputs are
content-hashed, whose seeds biject with the registered task seeds, and whose
group-bootstrap confidence intervals resample recordings rather than windows.

This article contributes:
\begin{enumerate}
\item a formalized specification of reproducible industrial ML experiments:
an eight-state lifecycle with guarded transitions and six integrity
invariants (\S\ref{sec:spec});
\item an end-to-end extensible architecture that enforces the specification
at its kernel boundary for heterogeneous datasets and both classical and deep
model backends (\S\ref{sec:arch}, \S\ref{sec:impl});
\item an automated integrity-verification and provenance stack---model
checking, oracle-based property tests, a twelve-entry fault-injection
catalogue, and hash-recheckable result bundles (\S\ref{sec:impl},
\S\ref{sec:eval});
\item an empirical evaluation across heterogeneous datasets (bearings,
a gearbox, acoustic and vibration modalities), independent plugin packaging,
and execution environments, including measurements of the framework's own
overhead (\S\ref{sec:eval}).
\end{enumerate}

We are deliberate about two boundaries. The formal model verifies the
\emph{abstract guard logic} of the design, not the Python implementation; the
implementation is covered by the runtime checks, property tests, and mutation
results, and we say so wherever the distinction matters. And the framework
has, so far, one author; evidence from third-party integrators is being
solicited but is not part of this article's claims.
